\chapter{Microservizi}

Nell'ingegneria del software, l'architettura a microservizi è un modello architettonico che organizza un'applicazione in un insieme di servizi a grana fine e liberamente accoppiati che comunicano attraverso protocolli leggeri. Questo modello è caratterizzato dalla capacità di sviluppare e distribuire servizi in modo indipendente, migliorando la modularità, la scalabilità e l'adattabilità. Tuttavia, introduce un'ulteriore complessità, in particolare nella gestione dei sistemi distribuiti e della comunicazione tra servizi, rendendo l'implementazione iniziale più impegnativa rispetto a un'architettura monolitica. \\ \\
Non esiste una definizione unica e universalmente condivisa di microservizi. Tuttavia, essi sono generalmente caratterizzati da un'attenzione particolare alla modularità, con ogni servizio progettato intorno a una specifica funzionalità aziendale. Questi servizi sono liberamente accoppiati, distribuibili in modo indipendente e spesso sviluppati e scalati separatamente, consentendo una maggiore flessibilità e agilità nella gestione di sistemi complessi. L'architettura a microservizi è strettamente associata a principi quali la progettazione guidata dal dominio, la decentralizzazione dei dati e della governance e la flessibilità di utilizzare tecnologie diverse per i singoli servizi, in modo da soddisfare al meglio i loro requisiti. 

\newpage
